\section{An intrinsic sample bound at every finite width}
\label{sec:critical-intrinsic-finite-width-learning}

The exceptional-mass theorem yields a finite-width learning statement in
terms of the target's actual prior mass. This formulation separates the
capacity question from the further question of estimating that mass from
a teacher implementation.

For the original unconditioned critical Gaussian prior $\Pi_N$, write
\[
 \mathcal A_{N,f}=\{w:f(x)z_w(x)>1/2\text{ for every }x\in\{-1,1\}^d\}.
\]
When $\Pi_N(\mathcal A_{N,f})>0$, define the intrinsic target cost
\[
 a_N(f)=-\frac1N\log\Pi_N(\mathcal A_{N,f}).
\]
This is the raw prior mass, before global-margin conditioning.
The learner conditions the original prior on an absolute output margin
at every input and on the exact training signs, as throughout the paper.

\begin{theorem}[Every-width learning from intrinsic target cost]
\label{thm:critical-intrinsic-finite-width-learning}
There are universal $c,C>0$ such that the following holds. Let $d\ge1$,
$L\ge2$, $N\ge1$, and suppose $\Pi_N(\mathcal A_{N,f})>0$.
Put $a=a_N(f)$. For any input distribution $D$ on the cube and
$0<\varepsilon,\delta<1$, a sufficient sample size is
\begin{equation}\label{eq:critical-intrinsic-finite-width-samples}
 n\ge\frac C\varepsilon\left\{
 \big[(d+L+a+2)(1+\log(1/\varepsilon))\big]^c
       \exp\!\big(C\sqrt{L(a-V_*)_+}\big)
       +\log(1/\delta)\right\}.
\end{equation}
With probability at least $1-\delta$ over the iid training inputs,
the posterior predictive probabilities satisfy
\[
 D\{x:p_N(f(x)\mid x)\le2/3\}\le\varepsilon.
\]
The theorem holds at the stated width $N$, with no sufficiently-large-width
qualification. Its dependence on $N$ is through the target's actual cost
$a_N(f)$.
\end{theorem}

\begin{proof}
Set $A=a+1/L$. By
Corollary~\ref{cor:critical-finite-width-exceptional-mass}, there is a
sample-independent class $\mathcal G_A$ with VC dimension at most $R$
and an exceptional-mass estimate valid for $N\ge R^{C_0}$, where, after
enlarging universal constants, we may use
\[
 R=(d+L+a+2)^{C_1}
            \exp\!\big(C_1\sqrt{L(a-V_*)_+}\big)\ge2.
\]
Indeed
$\sqrt{L(A-V_*)_+}\le\sqrt{L(a-V_*)_+}+1$.
Sauer's lemma on the $2^d$ cube inputs, or the trivial count when the
VC dimension exceeds the domain size, gives
\begin{equation}\label{eq:critical-finite-width-class-count}
 \log|\mathcal G_A|\le C(d+1)R.
\end{equation}

First suppose
\begin{equation}\label{eq:critical-intrinsic-large-width}
 N\ge R^{C_0},\qquad
 N\ge L\{dR\log2+\log(6/\varepsilon)\}.
\end{equation}
For every training sample, the raw mass of the posterior-conditioning
event is at least $\Pi_N(\mathcal A_{N,f})=e^{-Na}$.
The posterior mass outside $\mathcal G_A$ is consequently at most
\[
 \frac{\exp(dR\log2-NA)}{e^{-Na}}
   =\exp(dR\log2-N/L)\le\varepsilon/6.
\]
For a fixed Boolean function $g$ with test error greater than
$\varepsilon/6$, its chance of fitting all $n$ iid training labels is
at most $e^{-n\varepsilon/6}$. A union bound using
\eqref{eq:critical-finite-width-class-count} shows that, with probability
at least $1-\delta$, every consistent member of $\mathcal G_A$ has
error at most $\varepsilon/6$, provided
\[
 n\ge\frac6\varepsilon
       \{C(d+1)R+\log(1/\delta)\}.
\]
The posterior expected test error is then at most $\varepsilon/3$,
including its exceptional mass. On an input where the correct-label
posterior probability is at most $2/3$, the posterior error probability
is at least $1/3$. Averaging therefore proves the asserted confidence
criterion.

It remains to consider widths that do not satisfy
\eqref{eq:critical-intrinsic-large-width}. For a universal $C_2$,
every such width satisfies
\[
 N\le R^{C_2}\{1+\log(1/\varepsilon)\}.
\]
The VC dimension of all depth-$L$, width-$N$ tanh sign networks is at
most a universal polynomial in $d,L,N$. Applying the finite-cube count
once more, the logarithm of the number of Boolean sign functions in
the entire posterior support is at most
\[
 \operatorname{poly}(d,L,N)
 \le R^{C_3}\{1+\log(1/\varepsilon)\}^{C_3}.
\]
A union bound over this whole support shows that every consistent
function has test error at most $\varepsilon/6$ once
\eqref{eq:critical-intrinsic-finite-width-samples} holds with sufficiently
large universal constants. This implies the desired posterior guarantee
as above. Substituting the displayed formula for $R$ covers both width
ranges and proves the theorem.
\end{proof}

\paragraph{Scope of the finite-width conclusion.}
A bound on the limiting rate $J_L(f)$ alone does not give the same bound
on $a_N(f)$ at a specified finite width. The theorem is quantitative in
the actual raw target cost. The following sections supply separate
finite-width estimates of that cost for Fourier and circuit constructions. For example, if a separate
argument gives $a_N(f)\le V_*+K\log^2 Q/L$ and $d+L+a_N(f)\le Q^K$
for a universal $K$, then the displayed sample bound is polynomial in
$Q$, $1/\varepsilon$, and $\log(1/\delta)$ at that width. The target-based conclusions in
Section~\ref{sec:critical-all-width-target-learning} combine those
estimates with a VC bound for smaller widths.
